\documentclass[10pt,a4paper]{amsart}
\usepackage{amsmath,amssymb,amsthm,graphicx}
\usepackage{amsmath}
\usepackage{amssymb}
\usepackage{hyperref}
\renewcommand{\theenumi}{\Alph{enumi}}
\renewcommand{\thesubsection}{\thesection\Alph{subsection}}

%%%%%%%%%%%%%%%%%%%%%%%%%%%%%%%%%%%%%%%%%%
% ENVIRONMENTS
%%%%%%%%%%%%%%%%%%%%%%%%%%%%%%%%%%%%%%%%%%%
\newtheorem{theorem}{Theorem}[section]
\newtheorem{prop}[theorem]{Proposition}
\newtheorem{lemma}[theorem]{Lemma}
\newtheorem{remark}[theorem]{Remark}
\newtheorem{corollary}[theorem]{Corollary}
\newtheorem{definition}[theorem]{Definition}
\newtheorem{example}[theorem]{Example}
%
\def\beq{\begin{equation}}
\def\eeq{\end{equation}}
\def\be{\begin{equation*}}
\def\ee{\end{equation*}}
%%%%%%%%%%%%%%%%%%%%%%%%%
% Beginning of the text %
%%%%%%%%%%%%%%%%%%%%%%%%%
%\begin{document}



%\documentclass[10pt]{article}
\usepackage{amssymb}
\usepackage{amsfonts}
%\topmargin 0.3in
%\textwidth 6in
%\textheight 9in%\baselineskip
%\hoffset=-50pt
%\voffset=-2pt
%\newenvironment{proof}[1][Proof]
%{\ \rule{0.5em}{0.5em}}


\title{Dedekind-finite Cardinals having countable partitions}
\author{Supakun Panasawatwong and J K Truss} 
\date{This paper is an amplified version of part of the first author's PhD thesis at the University of Leeds}%\footnotemark}
\title{Every finitely branching level tree carries a sequence of ordered partitions with uniform successor numbers}
\author{Supakun Panasawatwong, J. K. Truss}
\date{arxiv:2204.05700v1; CC BY 4.0}
\begin{document}
\maketitle
\noindent\textit{Editorial note.} All delivered mathematical text is the author's own, copied from the pinned publisher source sources/dedekind-finite-trees-2204.05700v1/source0.tex; no statement or proof step was rewritten, reordered or removed.\par\medskip
\noindent\textit{Editorial note.} Preamble provenance: the retained preamble is the source's own layout adapter, if any, plus the author's macro block at lines 5-44 of sources/dedekind-finite-trees-2204.05700v1/source0.tex (everything between the publisher class metadata and \textbackslash{}begin\{document\}, comments excluded). Nothing was added, no macro was invented, and the publisher document class is deliberately not reproduced; the excerpt serves only as preamble, never as a span coordinate.\par\medskip
\noindent\textit{Editorial note.} Metadata provenance: source\_id and source\_version were taken from focused-seeds/logic/STK-logic-template-partitions/license-reuse.json; the title from the card spec, taken from the publisher abstract page meta; the author list from the card spec, taken from the publisher abstract page meta (citation\_author); the source URL from the first non-empty of source\_url, source\_urls[0] and the evidence's fixed\_version\_url/license\_url. The licence evidence file used by this card is focused-seeds/logic/STK-logic-template-partitions/license-reuse.json, recorded in license.evidence\_path so that the binding target is explicit.\par\medskip
\noindent\textit{Editorial note.} Licence binding: the licence evidence staged with this source (LICENSE\_EVIDENCE.json) carries no status field and none of the accepted notice forms, so this card points at a per-card reuse record written inside its own task directory (license-reuse.json) with the fields name / version\_url / prior\_license\_path / basis; it asserts a material basis only and makes no new licence claim.\par\medskip
\noindent\textit{Editorial note.} External references are navigation pointers only; every deliberately kept step of the author's argument stays inside the counted spans. The displayed numbers of those pointers are the author's own numbering under the author's \textbackslash{}newtheorem declarations (the source has no subsections, so numbering is section.index), computed mechanically by the lane's scanner.\par\medskip
\noindent\textit{Editorial note.} The counted proof contains no reference to any labelled item and no bibliography citation, so the card needs neither a navigation stub nor a bibliography excerpt.\par\medskip
\noindent\textit{Editorial note.} This is the sixth card this lane authors from this source and the last remaining low-cost candidate there; it was taken to keep the publishable queue fed before moving on to longer proofs.\par\medskip
\noindent\textit{Allowed standard prerequisites (editorial).}\begin{itemize}
\item The levels L\_k of a tree with omega levels, all finite, and the successor set of a node. Set up in the paragraphs preceding the counted statement (ending at source line 525). Fixes the objects that are partitioned and the successor counts that have to become uniform.
\item A sequence of natural numbers in which every number occurs infinitely often, and the eventual constancy of a refining sequence of partitions of a finite set. Used inside the counted proof at lines 532 and 546-547. Supplies the schedule of refinements and the reason the refining sequences stabilise.
\end{itemize}\medskip
\begin{theorem} Let $(T, \le)$ be a tree with $\omega$ levels, all finite, with $k$th level $L_k$. 
Then there is a sequence of ordered partitions $((\pi_k, <): k \in \omega)$ such that $\pi_k$ is a partition
of $L_k$, and such that for any $X \in \pi_k$ and $Y \in \pi_{k+1}$ there is $n(X,Y) \in \omega$ such that for
each $x \in X$ the number of successors of $x$ lying in $Y$ is equal to $n(X,Y)$. \label{4.2}  \end{theorem}

\begin{proof} We use the same sequence $(k_n)$ as in the previous proof, in which each member of $\omega$ 
appears infinitely often, and for each $k$ we find a refining sequence of partitions 
$((\pi_k^n, <): n \in \omega)$ of $L_k$, which here is taken to mean that each member of $\pi_k^n$ is a 
union of a convex subset of $\pi_k^{n+1}$. We start with each $\pi_k^0$ 
being the trivial partition of $L_k$ into just one piece.

Now assume that $(\pi_k^n, <)$ has been chosen for each $k$. If $k \neq k_n$, then $\pi_k^{n+1} = \pi_k^n$.
So we just have to define $\pi_k^{n+1}$ (and a suitable ordering of it) in the case that $k = k_n$. For each
$x \in L_k$, and $X \in \pi^n_{k+1}$, consider the sequence of natural numbers $(|X \cap succ(x)|: X \in \pi^n_{k+1})$
of length the cardinality of $\pi^n_{k+1}$ and let $x \sim y$ if the sequences corresponding to $x$ and $y$ are equal.
We then take $\pi^{n+1}_k$ to be the least common refinement of $\pi^n_k$ and the partition determined by $\sim$. We
can (definably) linearly order the finite sequences of natural numbers and hence extend the ordering of $\pi^n_k$ to
$\pi^{n+1}_k \upharpoonright X$ for each $X \in \pi^n_k$ which will thus form a convex subset.

Since $(\pi^n_k: n \in \omega)$ is a refining sequence and $L_k$ is finite, it must be eventually constant, at $\pi_k$ 
say. The relationship between $\pi_k$ and $\pi_{k+1}$ may now be described. Clearly as the sequence has stabilized,
$\sim$ must be trivial on each member of $\pi_k$. So if $X \in \pi_k$, $(|Y \cap succ(x)|: Y \in \pi_{k+1})$ is the same
sequence for each $x \in X$. Since $|succ(x)|$ is the sum of this sequence, this means in particular that each $x \in X$ 
has the same number of successors. And for each $Y \in \pi_{k+1}$, each $x \in X$ has the same number of successors in $Y$,
which we write as $n(X,Y)$ (which is allowed to be 0).     \end{proof} 

\end{document}
